\chapter{Background}

\section{Supernumerary robotic limbs}

Supernumerary robotic limbs augment rather than restore: the wearer keeps
their own limbs and gains additional ones. \todo{CITE: an origin paper for
supernumerary robotic limbs -- Parietti and Asada, bracing and body support}
Reported applications span bracing and postural support, supernumerary
fingers, and general-purpose extra arms for shared-workspace tasks.
\todo{CITE: supernumerary fingers -- Prattichizzo et al.}
\todo{CITE: a recent survey or quantitative review of SRL design and
embodiment}

Two questions recur in this literature and both matter here. The first is how
the extra limbs are \emph{commanded}, given that the wearer's own limbs are
already occupied; reported approaches include body-mounted sensing of residual
degrees of freedom, foot and tongue interfaces, and inference from the
wearer's ongoing motion. \todo{CITE: exploiting body redundancy to control
supernumerary robotic limbs -- Lisini Baldi et al., IJRR 2025}
\todo{CITE: shared control of SRLs using mixed reality and mouth-and-tongue
interfaces, Biosensors 2025} The second is whether the wearer comes to
experience the limbs as part of their body, which is normally framed in terms
of body schema and ownership. \todo{CITE: body schema / embodiment of extra
limbs}

The configuration used in this project sidesteps the first question entirely
by placing the operator outside the wearer's body, and makes the second
inapplicable: the arms are never mechanically linked to the wearer's own
limbs, so ownership language does not apply to them.

\section{Teleoperation with a separate wearer}

The direct predecessor to this configuration is Fusion, which mounted a
robotic head and two arms on a backpack worn by one person and drove them from
a second, remote operator wearing a head-mounted display.
\todo{CITE: Saraiji, Sasaki, Matsumura, Minamizawa and Inami, Fusion: full
body surrogacy for collaborative communication, ACM SIGGRAPH 2018 Emerging
Technologies, Article 7} Fusion defined three levels of what it termed bodily
driven communication, distinguished by how much the robot arms were physically
coupled to the wearer's own arms; in all three the operator commanded the
motion and the machine decided nothing. It was presented as an Emerging
Technologies demonstration and reported no user study or quantitative
evaluation, so it establishes the architecture as buildable rather than
characterised.

More recently, the spatial behaviour of back-mounted arms around a wearer has
been studied directly, with a concealed separate operator and varying levels
of apparent autonomy. \todo{CITE: SRL Proxemics: Spatial Guidelines for
Supernumerary Robotic Limbs in Near-Body Interactions, CHI 2026} That work
reported that higher autonomy did not improve the wearer's perceived safety,
and produced body-zone guidance on approach distances and trajectory shapes.
It is relevant to the present system both as a precedent for the two-person
arrangement and as a caution against assuming autonomy is experienced as
reassuring.

\section{Teleoperation from a moving base}

Where the manipulator base is not inertial, base motion enters the control
problem as a disturbance. Approaches reported for mobile manipulators include
disturbance estimation and compensation, and master--slave schemes with local
feedback on the follower.
\todo{CITE: motion/force control of mobile manipulators via extended
uncertainty and disturbance estimation}
For supernumerary limbs specifically, the wearer's body is the floating base
and the disturbance is the wearer's own postural sway.
\todo{CITE: Zhang et al. 2024, motion-compensation control of SRAs subject to
human-induced disturbances}
\todo{CITE: A Human Motion Compensation Framework for a Supernumerary Robotic
Arm, arXiv:2310.10029}
In the reported cases the wearer is also the operator, so the disturbance is
self-generated; the two-person configuration removes that, which is the
distinction this project is positioned against.

\section{Shared autonomy in teleoperation}

Assistive teleoperation is commonly formulated as an arbitration between the
operator's input and a robot policy derived from a prediction of their goal.
\todo{CITE: Dragan and Srinivasa, formalising assistive teleoperation /
policy blending}
A structural weakness of predicting before acting is that confidence often
arrives too late to be useful, which motivates formulations that assist under
goal uncertainty.
\todo{CITE: Javdani et al., shared autonomy via hindsight optimisation}
\todo{CITE: evidence on when blending helps and when it hurts}

\section{Platform: Kinova Gen3, ROS 2 and MoveIt 2}

The arms used are seven-degree-of-freedom Kinova Gen3 manipulators with
Robotiq 2F-85 grippers. \todo{CITE: Kinova Gen3 technical specification}
The redundant seventh degree of freedom is used in this work for collision
avoidance: the elbow can be swung about the shoulder--wrist axis without
moving the gripper, which provides a family of solutions to search when a pose
is otherwise blocked.

The software stack is ROS 2 Jazzy with MoveIt 2 for kinematics and collision
checking. \todo{CITE: ROS 2} \todo{CITE: MoveIt 2} Inverse kinematics is
solved by TRAC-IK. \todo{CITE: Beeson and Ames, TRAC-IK} TRAC-IK's use of
randomised restarts is relevant to the verification methodology described in
Chapter~\ref{ch:method}, because it makes a single successful solve a weak
statement about reachability.

\section{Open-vocabulary visual grounding}

The autonomy modes in this system require objects to be located from a
description rather than from a fixed model. Open-vocabulary detectors admit an
arbitrary class list at inference time.
\todo{CITE: YOLO-World open-vocabulary detection}
\todo{CITE: Grounding DINO}
Speech transcription uses a Whisper-family model.
\todo{CITE: Whisper / faster-whisper}
A deliberate decision not to use an end-to-end vision--language--action model
is discussed in Chapter~\ref{ch:discussion}.
\todo{CITE: OpenVLA or a comparable VLA, for the comparison being declined}
